\section{Virtualização de Redes (NV)}

A Virtualização de Redes (NV, \textit{Network Virtualization}) permite que múltiplas redes virtuais coexistam sobre uma mesma infraestrutura física, possibilitando o compartilhamento dos recursos disponíveis. Essa abordagem proporciona maior flexibilidade à infraestrutura de rede e permite que redes virtuais com características distintas sejam implementadas de forma independente sobre os mesmos recursos físicos \cite{chowdhury2010survey}.

A separação entre a infraestrutura física e as redes virtuais permite que diferentes configurações sejam estabelecidas sem a necessidade de modificar diretamente a estrutura física da rede. Dessa forma, cada rede virtual pode possuir sua própria organização e utilizar apenas uma parcela dos recursos disponibilizados pela infraestrutura compartilhada. Além de favorecer o aproveitamento desses recursos, essa independência proporciona maior flexibilidade para o gerenciamento e a configuração das redes virtuais \cite{chowdhury2010survey}.

Essa capacidade de compartilhar uma mesma infraestrutura física entre diferentes redes virtuais constitui uma característica importante para o fatiamento de redes. Por meio da virtualização, os recursos físicos podem ser abstraídos e utilizados por redes virtuais com diferentes configurações e requisitos. Em conjunto com outras abordagens, como SDN e NFV, a virtualização de redes fornece mecanismos que contribuem para a criação de fatias de rede customizáveis \cite{gomes2020strategies}.